%% =============================================================================
%% sm_e1_model_hierarchy.tex — SM-E, Subsection E.1: Model Hierarchy
%% Body content only (\subsection level); parent structure in sm_e.tex
%% =============================================================================

\subsection{Model hierarchy and architecture}
\label{sme:model-hierarchy}

\Cref{sme:tab-hierarchy} summarizes the five nested models fitted in
the application (\cref{sec:application}).  Each model extends its
predecessor by one structural feature: random intercepts (M1),
state-varying coefficients with block-diagonal covariance (M2),
cross-margin covariance (M3a), and policy moderators (M3b).  The
final model M3b-W adds survey weights to the pseudo-posterior.
Parameter counts include all fixed effects, random effects, and
hyperparameters (Cholesky factors, scales, and correlations).

The models are fitted sequentially using warm-start initialization:
the posterior mean of model~$k$ provides the initial values for
model~$k+1$.  Random-effect parameters absent in the simpler model
are initialized at zero (NCP variates) or identity (Cholesky factors).
This progressive strategy reduces total computation time by
approximately~$40\%$ compared with cold starts and ensures that each
model addresses a question raised by its predecessor:
M0 to M1 asks whether baseline enrolment rates vary by state;
M1 to M2, whether the poverty reversal varies by state;
M2 to M3a, whether cross-margin covariance improves fit;
and M3a to M3b, whether observable policies explain the state
heterogeneity.

\begin{table}[htbp]
  \centering
  \small
  \begin{tabular}{@{}l l r l r@{}}
    \toprule
    Model & Structure & Parameters
      & Covariance & $\widehat{\elpd}_{\mathrm{loo}}$ \\
    \midrule
    M0  & Pooled HBB
        & 11
        & ---
        & $-20{,}945$ \\
    M1  & Random intercepts
        & 116
        & $2 \times 2$
        & $-20{,}627$ \\
    M2  & Block-diagonal SVC
        & 551
        & $(5 \times 5) \oplus (5 \times 5)$
        & $-20{,}604$ \\
    M3a & Full SVC
        & 576
        & $10 \times 10$
        & $-20{,}602$ \\
    M3b & Policy moderators
        & 616
        & $10 \times 10$
        & $-20{,}608$ \\
    \bottomrule
  \end{tabular}
  \caption{Model hierarchy.  Parameters: total number of free
    parameters (fixed effects + random effects + hyperparameters).
    Covariance: dimension of the state-level covariance matrix for
    the random-effect deviations~$\bdelta_s$.  LOO: expected
    log-predictive density (\cref{sec:application}).  The
    progression from M0 to M3a yields significant incremental
    improvements in predictive accuracy; M3b adds policy
    moderators without further predictive gain.}
  \label{sme:tab-hierarchy}
\end{table}
